\documentclass[10pt,a4paper]{amsart}
\usepackage[margin=19mm]{geometry}
\usepackage{amsmath,amssymb,mathtools}
\usepackage[scr=rsfs,cal=euler]{mathalfa}
\usepackage{xfrac}
\usepackage[unicode,hidelinks,bookmarks=false]{hyperref}
\newcommand{\ie}{i.e.\ }
\newcommand{\cf}{cf.\ }
\newcommand{\resp}{respectively\ }
\newcommand{\Subsection}{\subsection}
\providecommand{\nobreakdash}{\nobreak\hbox{-}}
\setlength{\emergencystretch}{2em}
\multlinegap0pt
\usepackage{xfrac}
\usepackage{amssymb}
\usepackage{mathtools}
\usepackage{enumerate}
\newtheorem{lem}{Lemma}
\title{Tangential limits of stable minimal capillary surfaces}
\author{Michael Eichmair, Thomas Koerber}
\date{Source: arXiv:2605.12277v1 (CC BY 4.0)}
\begin{document}
\maketitle

\noindent\emph{Source and licence.} Tangential limits of stable minimal capillary surfaces. Version arXiv:2605.12277v1 (CC BY 4.0), distributed by arXiv under the Creative Commons Attribution 4.0 International licence, http://creativecommons.org/licenses/by/4.0/, as recorded in the arXiv licence metadata for this exact version and shown on the abstract page https://arxiv.org/abs/2605.12277v1. Full licence notice: \texttt{licenses/arxiv-2605.12277-CC-BY-4.0.txt}.\par\smallskip

\noindent\textit{Editorial scope.} Counted unit: the article's lem iteration (source0.tex lines 396--411). The article's complete original proof of it is delivered with it (source0.tex lines 412--520, one \texttt{proof} environment of 109 lines). No mathematics is written, repaired or re-derived: the statement and the proof are the article's own lines, in the article's own notation, and no retained line is substituted or re-flowed.  Retained uncounted as the source-local context the proof needs: lines 313--315; lines 322--324; lines 328--334; lines 336--341; lines 345--347; lines 2989--3001; lines 3107--3115; lines 3125--3137.  Nothing was omitted from any retained span, so the package carries no omission record.  No retained line contains a citation, so the wrapper carries no bibliography and no bibliography entry.  No retained line contains a figure environment, an image-inclusion command or drawing code, so no image is delivered and the package declares no asset dependency.  Editorial formatting only, in addition: an AMS \texttt{amsart} preamble that restates the article's own declarations and macros used by the retained lines, namely the environments \texttt{lem} on separate counters, together with the standard packages those lines need.  The retained environments print in this document as Lemma 1 in order of appearance, whereas the article prints the same items with the numbering of its own full document.  The complete native source of the article as distributed in the arXiv e-print for this exact version is bundled unchanged as \texttt{sources/200414/source0.tex}.
	\begin{align} \label{geometry bound} 
	1\leq \Lambda\,i_S(y)
	\end{align} 

	\begin{align} \label{H bound} 
	|H(S)(y)|\leq \Lambda\sin\theta
	\end{align} 

\begin{equation} \label{epsilon, delta} 
\begin{aligned}
	&\circ \qquad 2\,\varepsilon<r_2,\\
	&\circ\qquad 4\,\delta<\varepsilon,\text{ and}\\
	&\circ\qquad 4\,\delta\,\Lambda<1.
\end{aligned}
\end{equation} 

\begin{equation} \label{r0, r1}
	\begin{aligned} 
&\circ \qquad\varepsilon<r_1<r_2-\varepsilon\text{ and}\\
&\circ \qquad r_0<r_1-\varepsilon
	\end{aligned} 
\end{equation}

\begin{align} \label{assumption 0}
	\frac{1}{2}\,(-\tan\theta)\operatorname{dist}_\Sigma(x, \partial \Sigma)\leq \operatorname{dist}(x,S)\leq 2\,(-\tan\theta)\operatorname{dist}_\Sigma(x, \partial \Sigma)
\end{align}

	\begin{lem} \label{distance super harmonic}
		

 Let $S\subset \mathbb{R}^3$ be a support surface and  
$U\subset \mathbb{R}^3$  open. Assume that there is $\Lambda>1$ such that 
$$
1\leq \Lambda\,i_S(y)
$$
for all $y\in U\cap S$. Let $\Sigma\subset D(S)$ be a two-sided properly embedded minimal surface with  $U\cap \partial \Sigma=\emptyset$.  There are $w\in L^\infty(\Sigma)$ and $X\in L^\infty(\Sigma,\mathbb{R}^3)$ such that, on $$\{x\in \Sigma:\Pi_S(x)\in U\cap S\text{ and }2\,\Lambda\operatorname{dist}(x,S)<1\},$$ $|w|\leq 4\,\Lambda^2$, $|X|\leq 4\,\Lambda$,  and
		$$
		-\Delta^\Sigma \operatorname{dist}(\,\cdot\,,S)+X\cdot \nabla^\Sigma \operatorname{dist}(\,\cdot\,,S)+w\operatorname{dist}(\,\cdot\,,S)=H(S)\circ \Pi_S.
		$$ 
	\end{lem}

\begin{lem} \label{harnack inequality}
	Given $A>1$, there is $c>1$ with the following property. Let  $\Sigma\subset \mathbb{R}^3$ be a two-sided properly embedded minimal surface with $0\in  \Sigma$ and $r>0$ such that  $B^\Sigma_{2\,r}(0)\cap \partial \Sigma=\emptyset$ and  $B^\Sigma_{2\,r}(0)$ is stable.  Assume that $f\in C^\infty(B^\Sigma_{2\,r}(0))$ is positive and satisfies
	$$-\Delta^\Sigma f+X\cdot \nabla^\Sigma f+w\,f=q$$
	for some $w\in L^\infty(\Sigma)$ with $r^2\,|w|\leq A^2,$ $X\in L^\infty(\Sigma,\mathbb{R}^3)$  with $r\,|X|\leq A,$ and $q\in L^\infty(\Sigma)$.
	 Then 
	\begin{align}
		\sup\{f(x):x\in B^\Sigma_{r}(0)\}\leq c\,\inf \{f(x):x\in B^\Sigma_{r}(0)\}+c\,r^2\,\sup\{|q(x)|:x\in B^\Sigma_{2\,r}(0)\}.
	\end{align} 
\end{lem}

\begin{lem} \label{slab lemma}
Given $B>1$, there is $c>1$ such that, for all sufficiently small $0<\varepsilon<\sfrac{1}{2\,B}$, the following holds.  Let $r>0$, $S\subset \mathbb{R}^3$ be a support surface,   
 and $\Sigma\subset D(S)$ be a two-sided properly embedded minimal surface with $0\in \Sigma$ such that
 \begin{align}
 	&\circ\qquad r<B\,i_S(y)\text{ for all $y\in B_{2\,r}(0)\cap S$,} \notag \\
 	&\circ \qquad \text{$B^\Sigma_{2\,r}(0)\cap \partial \Sigma=\emptyset,$   $B^\Sigma_{2\,r}(0)$ is stable, and}\notag \\
 	&\circ\qquad  \operatorname{dist}(x,S)\leq r\,\varepsilon \text{  for all $x\in B^\Sigma_{2\,r}(0)$.}\notag 
 \end{align}
   There is   $u\in C^\infty(S)$ positive such that 
\begin{align*} 
r^{-1}\,|u|+|\nabla^S u|+r\,|\nabla^S\nabla^Su|\leq c\,\varepsilon
\end{align*} 
and $B^\Sigma_{r}(0)\subset \{y-u(y)\,\nu(S)(y):y\in S\}$.\end{lem} 

	\begin{lem}\label{harnack iteration}
Let   $L>1$.  There are $\sfrac{\pi}2<\theta_0<\pi$ and $c>1$ depending only on $\Lambda$, $\varepsilon$, $\delta$,  and $L$ such that the following holds. Let $\theta_0<\theta<\pi$ and
  $\alpha:[0,T]\to B_{r_1}(0)\cap\Sigma$ be a smooth curve of length at most $L$ with $\alpha(0)\in \partial \Sigma$. For all $t\in[0,T]$, there holds $$2\operatorname{dist}(\alpha(t),S)<\min\{\varepsilon,\,i_S(\Pi_S(\alpha(t)))\}$$ and 
\begin{align} \label{conclusion} 
\operatorname{dist}(\alpha(t),S)\leq c\,(-\tan\theta)\,s(t)
\end{align} 
where $s(t)=\min\{\operatorname{dist}_\Sigma(\alpha(t),\partial \Sigma),\sfrac{\delta}2\}$.
Moreover, for each $t\in[0,T]$, there is  $u\in C^\infty(S)$ positive such that
\begin{align} \label{graph concl 1}
\left\{z\in \Sigma:2\operatorname{dist}_\Sigma(\alpha(t),z)\leq s(t)\right\}\subset \{\Phi_S(u)(y):y\in S\}
\end{align} and
\begin{align} \label{graph concl2}
	|\nabla^Su(y)|+s(t)\,|\nabla^S\nabla^Su(y)|\leq c\,(-\tan\theta)
\end{align}
for all $y\in S$ with $\Phi_S(u)(y)\in \left\{z\in \Sigma:4\operatorname{dist}_\Sigma(\alpha(t),z)\leq s(t)\right\}$.
	\end{lem} 

	\begin{proof}
By \eqref{geometry bound} and \eqref{epsilon, delta},		 for all $y\in B_{r_2}(0)\cap S$,
		\begin{align} \label{delta 0}
			4\,\delta < \min\{\varepsilon,\,i_S(y)\}.
		\end{align}
		
		Let  $c_1>1$ be the constant from Lemma \ref{harnack inequality} applied with $A=\Lambda\,\delta$. 	Let  $c_2>1$ be the constant from Lemma \ref{slab lemma} applied with $B=\sfrac{\delta\,\Lambda}{4}$. Let $0<\varepsilon_2<1$ be such that Lemma \ref{slab lemma} holds for all $\varepsilon\in(0,\varepsilon_2)$ for this choice of $B$. Let $c_3=(1+\sfrac{\delta\,\Lambda}{4})\,c_1$.
		
		We choose  $\sfrac{\pi}2<\theta_0<\pi$ such that 
				\begin{align}  
			\label{theta small 0}
			4\,(-\tan\theta_0)\,(2+c_3^{\sfrac{4\,L}{\delta}})<\varepsilon_2.		
		\end{align}	
		
		It suffices to prove the assertion of the lemma for $t=T$. We may and will assume that $\alpha(T)\notin \partial \Sigma$.
			
		Assume first that $2\operatorname{dist}_\Sigma(\alpha(T),\partial \Sigma)< \delta$. 
		
		By \eqref{assumption 0}, \eqref{conclusion} holds with $c=2$.  Let $z\in \Sigma$ be such that
		$$
		\operatorname{dist}_\Sigma(\alpha(T),z)<\operatorname{dist}_\Sigma(\alpha(T),\partial \Sigma).
		$$ 
		 Note that $z\notin \partial \Sigma$. By \eqref{r0, r1} and \eqref{delta 0}, \begin{align} \label{z} 2\,|z|\leq 2\,|\alpha(T)|+2\,|\alpha(T)-z|<2\,r_1+\delta<2\, r_2.\end{align}
Thus, $z\in B_{r_2}(0)\cap \Sigma$. Moreover,	
		 $$
		 \operatorname{dist}_\Sigma(z,\partial \Sigma)\leq  \operatorname{dist}_\Sigma(z,\alpha(T))+ \operatorname{dist}_\Sigma(\alpha(T),\partial \Sigma)<2\operatorname{dist}_\Sigma(\alpha(T),\partial \Sigma)<\delta.
		 $$
		By \eqref{assumption 0} and \eqref{theta small 0}, 
		\begin{align} \label{close to boundary 0} 
		\operatorname{dist}(z,S)\leq 2\,(-\tan\theta)\operatorname{dist}_\Sigma(z,\partial \Sigma)<4\,(-\tan\theta)\operatorname{dist}_\Sigma(\alpha(T),\partial \Sigma)<2\,\delta.
		\end{align} 
		By \eqref{r0, r1}, \eqref{delta 0},  \eqref{theta small 0}, \eqref{z}, and \eqref{close to boundary 0}, 
		$$
		|\Pi_S(z)|\leq |z|+\operatorname{dist}(z,S)\leq r_1+3\,\delta<r_1+\varepsilon<r_2.
		$$
		Thus, $\Pi_S(z)\in B_{r_2}(0)\cap S$.
	By \eqref{delta 0}  and \eqref{close to boundary 0},  $2\operatorname{dist}(z,S)< \min\{\varepsilon,\,i_S(\Pi_S(z))\}$. By \eqref{geometry bound},
		$$
		\frac12\operatorname{dist}_\Sigma(\alpha(T),\partial \Sigma)<\frac{\delta}4=\frac{B}{\Lambda}\leq B\,i_S(y)
		$$
		for all $y\in B_{r_2}(0)\cap S$. By \eqref{theta small 0} and \eqref{close to boundary 0},  $$2\operatorname{dist}(z,S)\leq\varepsilon_2\operatorname{dist}_\Sigma(\alpha(T),\partial \Sigma).$$
		By Lemma \ref{slab lemma}, there is a  $u\in C^\infty(S)$ positive that satisfies \eqref{graph concl 1} and \eqref{graph concl2} with $c=8\,c_2$.
		
		
		
		Assume now that $2\operatorname{dist}_\Sigma(\alpha(T),\partial \Sigma)\geq \delta$.
		

By assumption,  the length of $\alpha:[0,T]\to B_{r_1}(0)\cap \Sigma$ is at most $L$. By \eqref{r0, r1} and \eqref{delta 0}, there are an integer $k\geq 1$ with 
		\begin{align} \label{k bound 0}
			(k+1)\,\delta\leq  4\,L
		\end{align}
		and  $0<t_1<\ldots <t_{k+1}=T$ such that 
		\begin{align} \label{p -1}
		\sfrac\delta4<	\operatorname{dist}_\Sigma(\alpha(t_1), \partial \Sigma)<\sfrac\delta2 		
		\end{align} 
	 and, for each $1\leq \ell \leq k$, 
		\begin{align}
			&\circ\qquad  \text{$\alpha(t_{\ell})\in B^\Sigma_{\sfrac{\delta}{4}}(\alpha(t_{\ell+1}))$,} \label{connect 0} \\
			&\circ \qquad \text{$B^\Sigma_{\sfrac{\delta}{2}}{(\alpha(t_{\ell+1}))}\cap \partial \Sigma=\emptyset$, and} \label{sufficiently large ball 1 0}	  \\
			&\circ \qquad \text{$B^\Sigma_{4\,\delta}{(\alpha(t_{\ell+1}))}\subset B_{r_2}(0)\cap\Sigma$.}	\label{sufficiently large ball 2 0}	 		
		\end{align}
		We claim that 
		\begin{align}
			&\circ \qquad \text{$2\operatorname{dist}(z,S)<\min\{\varepsilon,\,i_S(\Pi_S(z))\}$ for all  $z\in B^\Sigma_{\sfrac{\delta}{2}}(\alpha(t_\ell))$ and} \label{to show 1 0}\\
			&\circ\qquad  
			\operatorname{dist}(z,S)\leq c_3^{\ell-1}\,\delta\,(-\tan\theta)  \text{ for all  $z\in B^\Sigma_{\sfrac{\delta}{4}}{(\alpha(t_\ell))}$} \label{to show 2 0} 
		\end{align}			
		for each $2\leq \ell \leq k+1$.
		To see this, we proceed by induction. 
		
		First, if $\ell=2$ and $z\in B^\Sigma_{\sfrac{\delta}{2}}(\alpha(t_2))$, by  \eqref{assumption 0},  \eqref{theta small 0}, \eqref{p -1},  and \eqref{connect 0}, 
		$$
		\operatorname{dist}(z,S)\leq |z-\alpha(t_2)|+|\alpha(t_2)-\alpha(t_1)|+\operatorname{dist}(\alpha(t_1),S)<  	\frac{\delta}{2}+\frac{\delta}{4}	
		+(-\tan\theta)\,\delta<2\,\delta.
		$$
		In conjunction with \eqref{r0, r1} and \eqref{delta 0}, we have 
		$$
		|\Pi_S(z)|\leq|z|+\operatorname{dist}(z,S)\leq r_1+3\,\delta<r_1+\varepsilon<r_2.
		$$
		Using \eqref{delta 0} again, we obtain
		$$ 
			2\operatorname{dist}(z,S)< \min\{\varepsilon,\,i_S(\Pi_S(z))\}.
			$$
		In view of this, \eqref{H bound}, \eqref{connect 0}, \eqref{sufficiently large ball 1 0}, and \eqref{sufficiently large ball 2 0}, we may use Lemma \ref{distance super harmonic} and Lemma \ref{harnack inequality} with $f=\operatorname{dist}(\,\cdot\,,S)$ and $r=\sfrac{\delta}{4}$ to conclude that
			\begin{align*} 
			\operatorname{dist}(z,S)&\leq c_1\operatorname{dist}(\alpha(t_1),S)+\frac{\delta^2}{16}\,c_1\,\Lambda\,\sin\theta
			\end{align*} 
		 for all $z\in B^\Sigma_{\sfrac{\delta}{4}}(\alpha(t_2))$.  By \eqref{assumption 0} and \eqref{p -1}, 
		\begin{align*} 
	\operatorname{dist}(z,S)\leq c_1\,\delta\,(-\tan\theta)+\frac{\delta^2}{16}\,c_1\,\Lambda\,\sin\theta
	< c_3\,\delta\,(-\tan\theta).
		\end{align*} 
		Assume now that $2\leq \ell\leq k$ and that both \eqref{to show 1 0} and \eqref{to show 2 0} have already been shown for $\ell$. We obtain, for  $z\in B^\Sigma_{\sfrac{\delta}{2}}(\alpha(t_{\ell+1}))$, using  \eqref{delta 0}, \eqref{theta small 0}, \eqref{k bound 0}, \eqref{connect 0}, and \eqref{to show 2 0},  
		$$
		\operatorname{dist}(z,S)\leq |z-\alpha(t_{\ell+1})|+|\alpha(t_{\ell+1})-\alpha(t_{\ell})|+\operatorname{dist}(\alpha(t_{\ell}),S)\leq \frac\delta2+\frac\delta 4+c_3^{\ell-1}\,\delta\,(-\tan\theta)\leq 2\,\delta 
		$$
		As above, this implies that $|\Pi_S(z)|<r_2$ and $2\operatorname{dist}(z,S)\leq \min\{\varepsilon,\,i_S(\Pi_S(z))\}$.  
		We may use Lemma \ref{distance super harmonic}, Lemma \ref{harnack inequality}, and \eqref{to show 1 0} to show that \eqref{to show 2 0} also holds for $\ell+1$. 
		
		Specifying \eqref{to show 2 0} to the case where $\ell=k+1$ and using \eqref{k bound 0}, we obtain \eqref{conclusion} with $c=2\,c_3^{\sfrac{4\,L}{\delta}}$. Likewise, specifying \eqref{to show 2 0} to the case where $\ell=k+1$, using Lemma \ref{slab lemma}, \eqref{k bound 0}, and \eqref{theta small 0}, we see that there is  $u\in C^\infty(S)$ positive satisfying \eqref{graph concl 1} and \eqref{graph concl2} with $c=8\,c_2\,c_3^{\sfrac{4\,L}{\delta}}$.
		
	
		
		
		This completes the proof of the lemma.
		

	\end{proof}

\end{document}
